%% ma: L10-B15-0103
%% lop: 10
%% chuong: 6
%% bai: 15
%% dang: 10.15.7
%% loai: TN4
%% muc: NB
%% dap_an: "D"
%% nguon: "Ảnh giáo viên gửi 25/11/2026 (lô 2): Bài tập trắc nghiệm – Toán 10 Bài 15 Hàm số (sự biến thiên, đồ thị), Phần I Câu 14"
%% kien_thuc: "Lớp 10 Bài 15 (đồ thị của hàm số)"
%% trang_thai: chinh-thuc
%% ghi_chu: "Phương án C gốc là y=|x| với x<0. Hình gốc vẽ đến O; bản kho để hở điểm O cho khớp điều kiện x<0."
\begin{ex}
Đồ thị hình vẽ là đồ thị của một hàm số trong bốn hàm số được liệt kê ở bốn phương án A, B, C, D dưới đây. Hỏi hàm số đó là hàm số nào?
\begin{center}\begin{tikzpicture}[scale=.6]
\draw[truc] (-4.5,0)--(1.5,0) node[below]{$x$};
\draw[truc] (0,-.6)--(0,5) node[left]{$y$};
\draw[phu] (-1,0)--(-1,1)--(0,1);
\node[below] at (-1,0) {\footnotesize$-1$}; \node[below right] at (0,0) {\footnotesize$O$};
\node[right] at (0,1) {\footnotesize$1$};
\draw[thick] (-4.3,4.3)--(0,0);
\draw[fill=white] (0,0) circle (2pt);
\end{tikzpicture}\end{center}
\choice{$y=|x|$}{$y=-x$}{$y=x$ với $x<0$}{\True $y=-x$ với $x<0$}
\loigiai{Đồ thị chỉ nằm bên trái trục tung và đi qua $(-1;1)$, ứng với hàm số $y=-x$ với $x<0$.}
\end{ex}
